\documentclass[11pt,a4paper]{article}

% Packages
\usepackage[utf8]{inputenc}
\usepackage[T1]{fontenc}
\usepackage{amsmath, amssymb, amsthm}
\usepackage{graphicx}
\usepackage{booktabs}
\usepackage{hyperref}
\usepackage{geometry}
\usepackage{enumitem}
\usepackage{fancyhdr}
\usepackage{titlesec}
\usepackage{parskip}
\usepackage{xcolor}

% Page geometry
\geometry{margin=2.5cm, headheight=14pt}

% Color definitions
\definecolor{ImperialBlue}{RGB}{0, 62, 116}
\definecolor{DarkGray}{RGB}{51, 51, 51}
\definecolor{AccentOrange}{RGB}{214, 96, 77}

% Hyperref settings
\hypersetup{
    colorlinks=true,
    linkcolor=ImperialBlue,
    urlcolor=ImperialBlue,
    citecolor=ImperialBlue
}

% Header and footer
\pagestyle{fancy}
\fancyhf{}
\fancyhead[L]{\textcolor{ImperialBlue}{\small Blockchain Economics and Digital Assets}}
\fancyhead[R]{\textcolor{ImperialBlue}{\small Lecture 6 Notes}}
\fancyfoot[C]{\thepage}
\renewcommand{\headrulewidth}{0.4pt}

% Section formatting
\titleformat{\section}
  {\Large\bfseries\color{ImperialBlue}}
  {\thesection}{1em}{}
\titleformat{\subsection}
  {\large\bfseries\color{ImperialBlue}}
  {\thesubsection}{1em}{}
\titleformat{\subsubsection}
  {\normalsize\bfseries\color{DarkGray}}
  {\thesubsubsection}{1em}{}

% Custom commands
\newcommand{\keyword}[1]{\textbf{\textcolor{AccentOrange}{#1}}}

%----------------------------------------------------------------------------------------
\begin{document}

\begin{center}
    {\LARGE\bfseries\textcolor{ImperialBlue}{Blockchain Economics and Digital Assets}}\\[0.5em]
    {\Large Lecture 6: Digital Ownership and Tokenization}\\[1em]
    {\large Dr Daniele Bianchi}\\[0.5em]
    {\normalsize Queen Mary, University of London}\\[0.3em]
    {\normalsize Semester B, 2025/2026}
\end{center}

\vspace{1em}
\hrule
\vspace{1em}

\tableofcontents

\vspace{2cm}

%----------------------------------------------------------------------------------------
\section*{Overview}
\addcontentsline{toc}{section}{Overview}
%----------------------------------------------------------------------------------------

Previous lectures focused on blockchain as infrastructure for payments, programmable contracts, and stable-value tokens. This lecture examines a different application: representing \textit{ownership} of assets as tokens on a blockchain.

\keyword{Tokenization} is the process of representing rights to an asset---digital art, real estate, securities, credentials---as tokens on a blockchain. The token becomes a verifiable, transferable representation of ownership that can be traded, fractionalized, and programmed.

The 2021 NFT boom brought tokenization to mainstream attention, with digital images selling for millions of dollars. The subsequent crash exposed the speculative excess: wash trading inflated volumes, utility promises proved empty, and most collections lost nearly all value. But beneath the hype, the underlying technology---provable digital ownership---remains valuable.

Today, institutional interest has shifted toward \keyword{real-world asset} (RWA) tokenization. BlackRock, Franklin Templeton, and JPMorgan are tokenizing Treasury bills, money market funds, and other traditional assets. This represents a potentially more durable application of the same underlying capability.

This lecture covers the technology (token standards, metadata storage), the NFT market's rise and fall (and lessons learned), real-world asset tokenization, security tokens, and practical applications from credentials to supply chain tracking.

%----------------------------------------------------------------------------------------
\section{What is Tokenization?}
%----------------------------------------------------------------------------------------

\subsection{The Core Idea}

\keyword{Tokenization} represents ownership of assets as tokens on a blockchain.

What can be tokenized?
\begin{itemize}
    \item Digital art and collectibles
    \item Financial securities (stocks, bonds)
    \item Real estate
    \item Commodities
    \item Credentials and identity
    \item Virtually any asset with defined ownership
\end{itemize}

Why put ownership on-chain?
\begin{itemize}
    \item \textbf{Programmability}: Automate dividends, royalties, compliance
    \item \textbf{24/7 trading}: No market hours, instant settlement
    \item \textbf{Fractional ownership}: Own 0.01\% of a building
    \item \textbf{Transparency}: Verifiable ownership history
    \item \textbf{Global access}: Anyone with a wallet can participate
\end{itemize}

\subsection{Fungible vs Non-Fungible}

\keyword{Fungible} tokens are identical and interchangeable. One unit can be swapped for any other unit of the same type without difference. Examples: Bitcoin, ETH, USDC, company shares.

\keyword{Non-fungible} tokens (NFTs) are unique and not interchangeable. Each token has distinct properties or identifiers. Examples: digital art, event tickets, property deeds, credentials.

The key question is \textit{substitutability}: can you swap one for another without caring which specific one you have?

\subsection{What Does a Token Actually Contain?}

\textbf{On-chain} (stored directly on blockchain):
\begin{itemize}
    \item Token ID (unique identifier)
    \item Owner address
    \item Contract address
    \item Transfer history
\end{itemize}

\textbf{Off-chain} (stored elsewhere, linked via URI):
\begin{itemize}
    \item Image/media file
    \item Description and attributes
    \item Additional metadata
\end{itemize}

\keyword{Critical point}: Most NFTs do not store the actual image on-chain---it is too expensive. They store a \textit{link} to the image. If the server hosting the image goes down, you own a broken link.

\subsection{Metadata Storage Options}

\begin{center}
\begin{tabular}{lll}
\toprule
\textbf{Storage} & \textbf{Pros} & \textbf{Cons} \\
\midrule
On-chain & Permanent, decentralised & Very expensive, size limits \\
IPFS & Decentralised, content-addressed & Needs pinning, can disappear \\
Centralised server & Cheap, flexible & Single point of failure \\
Arweave & Permanent, one-time payment & Cost, smaller ecosystem \\
\bottomrule
\end{tabular}
\end{center}

\textbf{IPFS} (InterPlanetary File System) is decentralised storage where files are addressed by their content hash, not location. The same content always produces the same address, but files only persist if someone ``pins'' them.

\textbf{Best practice}: Use IPFS or Arweave for important metadata. Many early NFT projects used centralised servers---some images are already lost.

%----------------------------------------------------------------------------------------
\section{Token Standards}
%----------------------------------------------------------------------------------------

\subsection{Why Standards Matter}

Token standards define how tokens are created, how ownership is tracked, how transfers work, and what functions wallets and marketplaces can expect.

\keyword{Interoperability} is the key benefit. A standard means any wallet or marketplace can interact with any token following that standard, without custom integration.

\subsection{ERC-20: Fungible Tokens}

The standard for fungible tokens on Ethereum (2015).

Key functions:
\begin{itemize}
    \item \texttt{totalSupply()}: How many tokens exist?
    \item \texttt{balanceOf(address)}: How many does this address own?
    \item \texttt{transfer(to, amount)}: Send tokens to another address
    \item \texttt{approve(spender, amount)}: Allow another contract to spend your tokens
\end{itemize}

Examples: USDC, USDT, UNI, LINK, most DeFi tokens.

ERC-20 enabled the 2017 ICO boom and DeFi composability---standardised tokens that any protocol could integrate.

\subsection{ERC-721: Non-Fungible Tokens}

The standard for NFTs on Ethereum (2018).

Key differences from ERC-20:
\begin{itemize}
    \item Each token has a unique \texttt{tokenId}
    \item \texttt{ownerOf(tokenId)}: Who owns this specific token?
    \item \texttt{tokenURI(tokenId)}: Where is this token's metadata?
    \item Tokens are indivisible---you cannot transfer 0.5 of an NFT
\end{itemize}

Examples: CryptoPunks, Bored Apes, most digital art NFTs.

Limitation: Each token requires a separate transaction to mint. Creating 10,000 NFTs means 10,000 transactions---expensive.

\subsection{ERC-1155: Multi-Token Standard}

A more efficient standard supporting both fungible and non-fungible tokens in one contract (2019).

Key advantages:
\begin{itemize}
    \item Batch transfers: Send multiple token types in one transaction
    \item Mixed fungibility: Some IDs can be fungible, others unique
    \item Gas efficient: Significantly cheaper for games and large collections
\end{itemize}

\textbf{Use case}: Gaming. Token ID 1 might be gold coins (fungible, 1 million exist). Token ID 2 might be health potions (fungible, 10,000 exist). Token ID 3 might be a legendary sword (non-fungible, only 1 exists). All managed in one contract.

\subsection{Comparison}

\begin{center}
\begin{tabular}{lccc}
\toprule
& \textbf{ERC-20} & \textbf{ERC-721} & \textbf{ERC-1155} \\
\midrule
Fungibility & Fungible & Non-fungible & Both \\
Divisible & Yes & No & Configurable \\
Batch transfer & No & No & Yes \\
Gas efficiency & Medium & Low & High \\
Use case & Currencies, shares & Art, collectibles & Gaming, mixed \\
\bottomrule
\end{tabular}
\end{center}

%----------------------------------------------------------------------------------------
\section{The NFT Market: Rise and Fall}
%----------------------------------------------------------------------------------------

\subsection{The 2021 Boom}

NFT trading volume exploded in 2021--2022, reaching billions of dollars per month. Bored Ape Yacht Club floor prices exceeded \$400,000. Beeple sold digital art for \$69 million at Christie's.

\textbf{What drove the boom}:
\begin{itemize}
    \item Crypto bull market creating wealth effect
    \item Social signalling (profile picture culture)
    \item Speculation and ``greater fool'' dynamics
    \item FOMO and media hype
    \item Genuine novelty of provable digital ownership
\end{itemize}

\subsection{The 2023 Crash}

Trading volume collapsed by over 90\%. Most collections lost 90\%+ of value. The ``blue chip'' collections (Bored Apes, CryptoPunks) retained some value; most others became worthless.

\textbf{What drove the crash}:
\begin{itemize}
    \item Crypto winter reduced speculative capital
    \item Utility promises (games, metaverse access) did not materialise
    \item Wash trading exposed---actual volumes were far lower than reported
    \item Oversupply: Too many collections, not enough genuine collectors
    \item Royalty enforcement failed---platforms competed by removing creator royalties
\end{itemize}

\subsection{The Wash Trading Problem}

\keyword{Wash trading} is buying and selling an asset to yourself (or colluding parties) to create fake trading volume.

Why it happened in NFTs:
\begin{itemize}
    \item Marketplace rewards (tokens, airdrops) based on volume
    \item Create appearance of liquidity and demand
    \item Manipulate ``floor price'' for collections
    \item Tax loss harvesting (in some jurisdictions)
\end{itemize}

Studies estimate 40--80\% of NFT volume was wash trading at peak. Detection methods: same wallet on both sides, circular transactions, economically irrational patterns, wallets funded from same source.

\subsection{NFT Valuation Challenges}

How do you value a unique asset with no cash flows?

Traditional valuation does not apply:
\begin{itemize}
    \item No dividends $\rightarrow$ No DCF
    \item No earnings $\rightarrow$ No P/E ratio
    \item Unique items $\rightarrow$ No exact comparables
\end{itemize}

What drives NFT prices:
\begin{itemize}
    \item Rarity of traits (verifiable on-chain)
    \item Artist/creator reputation
    \item Collection brand and community
    \item Historical sales of similar items
    \item Social signalling value
    \item Speculation on future demand
\end{itemize}

Prices are highly subjective and volatile---more art market than financial market.

\subsection{Lessons from the Bubble}

\textbf{What went wrong}:
\begin{itemize}
    \item Speculation overwhelmed utility
    \item ``Community'' and ``roadmap'' promises were often empty
    \item Royalties were not technically enforceable
    \item Legal status of ownership remained unclear
    \item Most buyers were speculators, not collectors
\end{itemize}

\textbf{What remains valuable}:
\begin{itemize}
    \item The \textit{technology} of provable digital ownership
    \item Creator royalties as a \textit{concept} (even if enforcement failed)
    \item On-chain provenance and authenticity verification
    \item Programmable ownership rights
\end{itemize}

The infrastructure is useful. The 2021 use case (speculative JPEGs) was not sustainable.

%----------------------------------------------------------------------------------------
\section{Real-World Asset Tokenization}
%----------------------------------------------------------------------------------------

\subsection{The Institutional Pivot}

As NFT speculation collapsed, institutional interest shifted to tokenizing traditional assets: Treasury bills, money market funds, corporate bonds, real estate.

\textbf{Key players entering}: BlackRock (BUIDL fund), Franklin Templeton (on-chain money market fund), JPMorgan (Onyx platform), Goldman Sachs (Digital Asset Platform).

This represents a more conservative but potentially durable application of tokenization technology.

\subsection{Tokenized Treasuries}

\textbf{What it is}: US Treasury bills represented as tokens on a blockchain.

\textbf{How it works}:
\begin{enumerate}
    \item Issuer (e.g., BlackRock) buys T-bills
    \item Issues tokens representing shares in the T-bill portfolio
    \item Token holders receive yield (distributed on-chain)
    \item Tokens can be transferred 24/7
\end{enumerate}

\textbf{Why it matters}:
\begin{itemize}
    \item DeFi protocols can hold yield-bearing collateral
    \item 24/7 settlement (vs T+1 for traditional)
    \item Programmable: Auto-reinvest, use as collateral in smart contracts
    \item Global access to US government debt
\end{itemize}

Market size (2024): Approximately \$1.5 billion in tokenized treasuries.

\subsection{Case Study: BlackRock BUIDL}

\textbf{BUIDL} (BlackRock USD Institutional Digital Liquidity Fund) launched March 2024 on Ethereum.

Structure:
\begin{itemize}
    \item Invests in T-bills, repos, cash
    \item Each BUIDL token = \$1 (stable NAV)
    \item Daily accrued dividends (paid monthly)
    \item Tokenised by Securitize
\end{itemize}

Requirements:
\begin{itemize}
    \item Qualified purchasers only (\$5M+ investable assets)
    \item \$5 million minimum investment
    \item KYC/AML through Securitize
\end{itemize}

\textbf{Significance}: The world's largest asset manager tokenizing products on a public blockchain signals institutional legitimacy.

\subsection{Why Tokenize Real-World Assets?}

\begin{center}
\begin{tabular}{lcc}
\toprule
& \textbf{Traditional} & \textbf{Tokenized} \\
\midrule
Settlement & T+1 or T+2 & Near-instant \\
Trading hours & Market hours & 24/7/365 \\
Minimum investment & Often \$1,000+ & Can be \$1 \\
Fractional ownership & Limited & Native \\
Cross-border & Complex, expensive & Simplified \\
Transparency & Periodic reports & Real-time on-chain \\
Programmability & None & Smart contracts \\
\bottomrule
\end{tabular}
\end{center}

The promise: more efficient capital markets with lower friction, broader access, and automated compliance.

\subsection{RWA Tokenization Challenges}

\textbf{Legal and regulatory}:
\begin{itemize}
    \item What does the token legally represent?
    \item Securities law compliance
    \item Cross-jurisdictional recognition
    \item Bankruptcy treatment of token holders
\end{itemize}

\textbf{Technical}:
\begin{itemize}
    \item Oracle problem: Who verifies the off-chain asset exists?
    \item Custody of underlying assets
    \item Blockchain scalability and costs
    \item Interoperability between chains
\end{itemize}

\textbf{Market structure}:
\begin{itemize}
    \item Liquidity fragmentation across platforms
    \item Need for qualified custodians
    \item Integration with existing financial infrastructure
\end{itemize}

Tokenization does not eliminate counterparty risk---it shifts it to whoever holds and verifies the underlying asset.

%----------------------------------------------------------------------------------------
\section{Security Tokens}
%----------------------------------------------------------------------------------------

\subsection{Definition}

A \keyword{security token} represents ownership in a regulated security---equity, debt, or investment contract---and is subject to securities law.

Key distinction:
\begin{itemize}
    \item \textbf{Utility token}: Access to a product/service (may not be a security)
    \item \textbf{Security token}: Investment contract, ownership stake (definitely a security)
\end{itemize}

\textbf{The Howey Test} (US): Is it an ``investment of money in a common enterprise with expectation of profits from efforts of others''? If yes, it is a security.

\subsection{ICOs and Their Downfall}

\keyword{Initial Coin Offerings} (ICOs) raised over \$6 billion in 2017 by selling tokens directly to investors, often before products existed.

Problems:
\begin{itemize}
    \item Most ICO tokens were likely unregistered securities
    \item Variable-quality whitepapers, no investor protections
    \item Anyone could participate regardless of sophistication
\end{itemize}

The SEC cracked down starting with the 2017 DAO Report. Enforcement actions followed: Munchee (\$15M returned), Block.one/EOS (\$24M settlement), Telegram/TON (\$1.2B returned, project cancelled).

Result: The ICO market collapsed. Projects seeking compliant fundraising turned to \keyword{Security Token Offerings} (STOs).

\subsection{STOs vs ICOs}

\begin{center}
\begin{tabular}{lll}
\toprule
& \textbf{ICO (2017)} & \textbf{STO} \\
\midrule
Regulatory status & Often unclear & Registered or exempt \\
Investor requirements & Anyone & Often accredited only \\
Disclosure & Whitepaper (variable) & Prospectus/offering memo \\
Investor protection & Minimal & Securities law applies \\
Secondary trading & Unregulated exchanges & Licensed platforms \\
\bottomrule
\end{tabular}
\end{center}

US exemptions for STOs: Reg D (accredited investors), Reg A+ (mini-IPO up to \$75M), Reg S (non-US investors).

\subsection{Security Token Infrastructure}

Issuance platforms (Securitize, Polymath) provide:
\begin{itemize}
    \item KYC/AML verification of all token holders
    \item Transfer restrictions (only to verified wallets)
    \item Compliance with holding periods
    \item Cap table management
    \item Reporting and disclosure
\end{itemize}

Technically: whitelists of approved addresses, transfer functions that check compliance before executing, forced transfers for legal requirements, automated dividend distribution.

This is \textit{permissioned} tokenization---not the ``anyone can participate'' ethos of DeFi.

\subsection{Current State}

What has been tokenized: real estate, private company equity, fund shares, corporate bonds, revenue-sharing agreements.

Challenges:
\begin{itemize}
    \item Limited secondary market liquidity
    \item High compliance costs relative to deal size
    \item Fragmented global regulatory landscape
    \item Chicken-and-egg: need liquidity to attract issuers, need issuers to build liquidity
\end{itemize}

Security tokens are growing but remain niche. Institutional adoption (BlackRock, etc.) may catalyse broader growth.

%----------------------------------------------------------------------------------------
\section{Practical Applications}
%----------------------------------------------------------------------------------------

\subsection{Digital Art and Collectibles}

Despite the crash, the core use case---provable digital ownership for art---remains valid.

\textbf{What works}: Provenance tracking, direct artist-to-collector sales, global market access.

\textbf{What does not}: Speculation-driven pricing, copyright transfer (usually not included), the ``right-click save'' criticism (though this misunderstands ownership vs copying).

The sustainable use is authentication and provenance---galleries and auction houses integrating NFTs for verification, not speculative trading.

\subsection{Supply Chain and Provenance}

Track products from origin to consumer using tokens as identifiers.

Examples: Food traceability (Walmart + IBM Food Trust), pharmaceutical supply chain, conflict mineral tracking, luxury goods authentication (LVMH's Aura).

How it works: Each product gets a unique token/identifier; each handoff is recorded on-chain; consumers can verify full history.

\textbf{Limitation}: The oracle problem again. Blockchain cannot verify that the physical item matches the digital record. ``Garbage in, garbage out''---immutably recorded.

\subsection{Credentials and Soulbound Tokens}

\keyword{Soulbound tokens} (SBTs) are non-transferable tokens representing credentials, affiliations, or reputation.

Examples: University degrees, professional certifications, proof of attendance, reputation scores.

Unlike NFTs, SBTs cannot be sold or transferred---they are ``bound'' to the recipient's wallet, representing earned or verified attributes rather than tradeable assets.

Challenges: Privacy (credentials visible on public blockchain), recovery (what if you lose wallet access?), adoption (issuers must participate).

\subsection{Fractional Ownership}

Tokenize expensive assets and sell fractions.

Examples: Real estate (own 0.1\% of a building), art (shares in a Picasso), collectibles (fractional sports memorabilia).

Benefits: Lower minimum investment, diversification, liquidity for illiquid assets.

Challenges: Securities law (usually applies), governance (who decides to sell the underlying?), custody and insurance, thin secondary markets.

Reality: Fractionalization existed pre-blockchain (REITs, art funds). Tokenization adds programmability and potentially broader access, but regulatory complexity remains.

%----------------------------------------------------------------------------------------
\section{Challenges and Future}
%----------------------------------------------------------------------------------------

\subsection{The Oracle Problem (Again)}

Blockchain can only verify on-chain data. For tokenized real-world assets, someone must verify: Does the asset exist? Is it authentic? Who owns it legally?

This gap between physical and digital worlds requires trusted intermediaries---auditors, custodians, legal systems. Tokenization does not eliminate trust; it shifts it.

\subsection{Legal Status}

Key unresolved questions:
\begin{itemize}
    \item Does the token holder have legal title to the underlying?
    \item What happens in bankruptcy?
    \item Which jurisdiction's law applies?
    \item How are disputes resolved?
\end{itemize}

Current approaches: Token represents contractual claim (not direct ownership), SPV structures (token = shares in entity owning asset), legal wrappers with token as ``digital twin.''

Regulatory developments: MiCA (EU framework), UK Law Commission (recognising digital assets as property), Wyoming (DAO and digital asset laws).

\subsection{The Future}

\textbf{Near-term} (2--5 years):
\begin{itemize}
    \item Tokenized treasuries and money market funds scale
    \item More traditional asset managers enter
    \item Regulatory clarity improves
    \item Security token infrastructure matures
\end{itemize}

\textbf{Medium-term} (5--10 years):
\begin{itemize}
    \item Broader asset classes tokenized (private equity, real estate)
    \item Integration with traditional financial infrastructure
    \item CBDCs potentially interoperating with tokenized assets
    \item Credentials and identity (SBTs) gain adoption
\end{itemize}

\textbf{Key question}: Will tokenization happen on public blockchains (Ethereum) or private/permissioned systems? The answer affects decentralisation, access, and composability.

%----------------------------------------------------------------------------------------
\section{Summary and Looking Ahead}
%----------------------------------------------------------------------------------------

This lecture has covered digital ownership and tokenization. Key takeaways:

\textbf{Tokenization represents ownership as blockchain tokens.} Fungible (ERC-20) for interchangeable assets, non-fungible (ERC-721) for unique items, multi-token (ERC-1155) for mixed use cases.

\textbf{NFT speculation crashed, but the technology remains.} Wash trading, hype, and unsustainable economics drove the bubble; provable digital ownership remains valuable.

\textbf{Real-world asset tokenization is growing.} BlackRock, Franklin Templeton, and others are tokenizing treasuries and funds, bringing institutional credibility.

\textbf{Security tokens require compliance.} STOs differ from ICOs in regulatory status, investor requirements, and secondary market restrictions.

\textbf{Practical applications extend beyond speculation.} Credentials, supply chain, fractional ownership---though all face the oracle problem.

\textbf{Challenges remain.} Legal status of tokens, physical-digital bridge, interoperability, and liquidity fragmentation.

In the next lecture, we turn to cryptocurrency investment---market structure, valuation approaches, risk factors, and the emergence of regulated investment products like ETFs.

%----------------------------------------------------------------------------------------
\section*{Readings}
\addcontentsline{toc}{section}{Readings}
%----------------------------------------------------------------------------------------

\textbf{Required:}
\begin{itemize}
    \item Chalmers, D., Fisch, C., Matthews, R., Quinn, W., \& Sherrill, J. (2022). ``Beyond the Bubble: Will NFTs and Digital Proof of Ownership Empower Creative Industry Entrepreneurs?'' \textit{Journal of Business Venturing Insights}, 17, e00309.
\end{itemize}

\textbf{Supplementary:}
\begin{itemize}
    \item Nadini, M., et al. (2021). ``Mapping the NFT Revolution: Market Trends, Trade Networks, and Visual Features.'' \textit{Scientific Reports}, 11, 20902.
    \item World Economic Forum. (2024). ``Tokenization of Real-World Assets: A Path to More Efficient Capital Markets.''
    \item Weyl, E. G., Ohlhaver, P., \& Buterin, V. (2022). ``Decentralized Society: Finding Web3's Soul.'' On soulbound tokens.
\end{itemize}

%----------------------------------------------------------------------------------------

\end{document}